\documentclass[11pt]{article}
\usepackage[a4paper,margin=1in]{geometry}
\usepackage{fontspec}
\usepackage[boldfont=false]{xeCJK}
\setCJKmainfont{FandolSong-Regular.otf}
\usepackage{amsmath}
\usepackage{amssymb}
\usepackage{array}
\usepackage{booktabs}
\usepackage{graphicx}
\usepackage{adjustbox}
\usepackage{enumitem}
\setlist[itemize]{topsep=2pt,partopsep=0pt,parsep=0pt,itemsep=2pt,leftmargin=1.4em}
\setlist[enumerate]{topsep=2pt,partopsep=0pt,parsep=0pt,itemsep=2pt,leftmargin=1.6em}
\usepackage{parskip}
\usepackage{fvextra}
\fvset{breaklines=true,breakanywhere=true,fontsize=\small}
\setlength{\parindent}{0pt}

\begin{document}

\begin{center}
  {\LARGE\bfseries Selection and evolution}\\[0.35em]
  {\large A Level Biology}
\end{center}
\vspace{0.6em}

\section*{Variation}

\textbf{Variation} 变异 means the differences between individuals. It has three possible causes:

\begin{itemize}
\item \textbf{genetic} factors only — set by the \textbf{alleles} 等位基因 you inherit (for example human blood group).

\item \textbf{environmental factors} 环境因素 only — set by your surroundings (for example a scar, or the language you speak).

\item a \textbf{combination} of both — most features, such as height and body mass, depend on genes \textbf{and} on diet and lifestyle.

\end{itemize}

There are two patterns of variation:

\begin{itemize}
\item \textbf{discontinuous variation} 不连续变异 — clear, separate groups with nothing in between (for example blood group A, B, AB or O). It is usually controlled by one or a few \textbf{genes} 基因, with little effect from the environment.

\item \textbf{continuous variation} 连续变异 — a smooth range from one extreme to the other (for example height). It is controlled by many genes together, plus the environment.

\end{itemize}

To compare the means of two samples (for example the heights of plants in sun and in shade), you use the \textbf{t-test}, which tells you whether the difference is large enough to be real, or just due to chance.

\section*{Natural selection}

A \textbf{population} 种群 produces far more \textbf{offspring} 后代 than can survive, so the offspring must \textbf{compete} 竞争 for resources such as food and space. This is the "struggle for existence". The individuals that are best \textbf{adapted} 适应 are most likely to survive, \textbf{reproduce} 繁殖, and pass on their alleles to the next generation. Over many generations, the helpful alleles become more common in the population. This is \textbf{natural selection} 自然选择.

Environmental conditions can push selection in three ways:

\begin{itemize}
\item \textbf{stabilising selection} 稳定选择 favours the average and removes the extremes (the population stays the same).

\item \textbf{directional selection} 定向选择 favours one extreme, so the mean shifts that way.

\item \textbf{disruptive selection} 分裂选择 favours both extremes and removes the average.

\end{itemize}

Allele frequencies can also change in other ways:

\begin{itemize}
\item the \textbf{founder effect} 奠基者效应 — a few individuals start a new population, so they carry only some of the alleles of the original group.

\item \textbf{genetic drift} 遗传漂变 — in a small population, allele frequencies change by chance from generation to generation.

\item the \textbf{bottleneck effect} 瓶颈效应 — a sudden fall in population size leaves few survivors, reducing the variety of alleles.

\end{itemize}

\subsection*{Antibiotic resistance as natural selection}

A chance \textbf{mutation} 突变 makes a few bacteria \textbf{resistant} to an \textbf{antibiotic} 抗生素. When the antibiotic is used, the non-resistant bacteria die, but the resistant ones survive and reproduce. Over time the \textbf{resistance} 耐药性 spreads through the population. This is natural selection in action.

\section*{The Hardy–Weinberg principle}

The \textbf{Hardy–Weinberg principle} 哈迪-温伯格原理 lets you calculate the \textbf{allele frequencies} 等位基因频率 and \textbf{genotype} 基因型 frequencies in a population. It only holds true when there is a large population, mating is random, and there is no mutation, no migration and no natural selection.

\section*{Selective breeding (artificial selection)}

In \textbf{selective breeding} 选择育种, also called \textbf{artificial selection} 人工选择, humans (not nature) choose which organisms breed, so that useful features are passed on. Examples:

\begin{itemize}
\item breeding \textbf{disease resistance} 抗病性 into varieties of wheat and rice.

\item using \textbf{inbreeding} 近交 and \textbf{hybridisation} 杂交 (crossing different lines) to make vigorous, uniform maize.

\item breeding dairy cattle to improve their milk \textbf{yield} 产量.

\end{itemize}

\section*{Evolution}

\textbf{Evolution} 进化 is the slow formation of new \textbf{species} 物种 from earlier ones, as the \textbf{gene pool} 基因库 (all the alleles in a population) changes from generation to generation.

DNA \textbf{sequence} 序列 data can show how closely related two species are: the more similar their DNA sequences, the more recently they shared a common ancestor.

\textbf{Speciation} 物种形成 happens when two populations become genetically separated, so they can no longer breed together. This genetic \textbf{isolation} 隔离 can come about in two ways:

\begin{itemize}
\item \textbf{allopatric speciation} 异域物种形成 — the populations are kept apart by a \textbf{geographical separation} 地理隔离, such as a sea or a mountain range.

\item \textbf{sympatric speciation} 同域物种形成 — the populations live in the same area but are separated by differences in behaviour or way of life.

\end{itemize}


\end{document}
